\documentclass[a4paper,12pt]{article}
\usepackage[utf8]{inputenc}
\usepackage[T1]{fontenc}
\usepackage[ngerman]{babel}
\usepackage{geometry}
\usepackage{amsmath}
\usepackage{hyperref}
\usepackage{csquotes}
\geometry{margin=2.5cm}

\begin{document}

\section*{NSFnets (Navier-Stokes Flow Nets): Physics-Informed Neural Networks
for the Incompressible Navier-Stokes Equations}

\textbf{Autoren:} Xiaowei Jin, Shengze Cai, Hui Li, George Em Karniadakis \\
\textbf{Jahr:} 2021 \\
\textbf{Quelle:} \textit{Computer Methods in Applied Mechanics and Engineering}, 376, 113627

% -----------------------------------------------------------------------
\subsection*{Kernbeitrag und Motivation}

Jin et al.\ präsentieren NSFnets als spezialisierte Variante der
Physics-Informed Neural Networks (PINNs) für die inkompressiblen
Navier-Stokes-Gleichungen. Die zentrale Motivation der Arbeit ist dieselbe,
die auch für Surrogatmodelle in der Geophysik gilt: klassische numerische
Strömungssimulationen (CFD) sind zu rechenintensiv für umfangreiche
Parameterstudien, weshalb dateneffiziente neuronale Netze als schnelle
Ersatzmodelle gesucht werden. Der Unterschied zu rein datengetriebenen
Ansätzen liegt darin, dass PINNs die Physik der Strömung explizit in die
Verlustfunktion integrieren -- sie lernen also nicht nur aus Beispieldaten,
sondern werden gleichzeitig durch die Differentialgleichungen selbst
regularisiert.

% -----------------------------------------------------------------------
\subsection*{Zwei Formulierungen der Navier-Stokes-Gleichungen}

Ein zentraler methodischer Beitrag der Arbeit ist der systematische Vergleich
zweier äquivalenter mathematischer Formulierungen der inkompressiblen
Navier-Stokes-Gleichungen als Grundlage für das Training:

\paragraph{Geschwindigkeit-Druck-Formulierung (VP).}
Das neuronale Netz bildet Raum-Zeit-Koordinaten $(t, x, y, z)$ auf die
Felder $(u, v, w, p)$ ab. Die Verlustfunktion bestraft Residuen der
Impulsgleichungen sowie der Inkompressibilitätsbedingung
$\nabla \cdot \mathbf{u} = 0$ als separate Terme. Druck ist dabei eine
\emph{versteckte Zustandsgröße}: Es werden keine Druckrandbedingungen
benötigt, da der Druck allein über die Divergenzfreiheitsbedingung implizit
bestimmt wird -- ohne die in klassischen CFD-Methoden übliche
Druck-Poisson-Gleichung.

\paragraph{Wirbelstärke-Geschwindigkeit-Formulierung (VV).}
Das Netz gibt statt Druck die Wirbelstärke $\boldsymbol{\omega} = \nabla
\times \mathbf{u}$ aus. Da $\nabla \cdot (\nabla \times \mathbf{u}) \equiv
0$ per Konstruktion gilt, ist die Divergenzfreiheitsbedingung bei dieser
Formulierung \emph{mathematisch automatisch erfüllt} -- die
Inkompressibilität muss nicht als Soft-Constraint in der Verlustfunktion
erzwungen werden. Die Verlustfunktion enthält stattdessen die
Wirbelstärke-Transport- sowie die Poisson-Gleichung für die Geschwindigkeit.
Die VV-Formulierung erwies sich für laminare Strömungen als präziser als
die VP-Variante, da ein expliziter Freiheitsgrad (Druck) entfällt und die
physikalische Nebenbedingung strukturell eingebettet ist.

% -----------------------------------------------------------------------
\subsection*{Verlustfunktion und dynamische Gewichtung}

Die gesamte Verlustfunktion für die VP-Formulierung lautet:
\begin{equation}
    \mathcal{L} = \mathcal{L}_e + \alpha\,\mathcal{L}_b + \beta\,\mathcal{L}_i,
\end{equation}
wobei $\mathcal{L}_e$ den physikalischen Residualterm (Navier-Stokes-Gleichungen),
$\mathcal{L}_b$ den Randbedingungs- und $\mathcal{L}_i$ den
Anfangsbedingungsterm bezeichnet. Die Gewichte $\alpha$ und $\beta$ steuern das
Verhältnis zwischen datengetriebenen und physikbasierten Anteilen.

Ein wesentlicher praktischer Beitrag ist die Einführung einer \emph{dynamischen
Gewichtungsstrategie}: Statt $\alpha$ und $\beta$ manuell einzustellen, werden
sie während des Trainings adaptiv aktualisiert, indem die maximalen
Gradienten des Physik-Residualterms ins Verhältnis zu den Gradienten der
Datenterme gesetzt werden:
\begin{equation}
    \hat{\alpha}^{(k+1)} = \frac{\max_\theta|\nabla_\theta \mathcal{L}_e|}
    {|\nabla_\theta \alpha^{(k)} \mathcal{L}_b|}.
\end{equation}
Die Koeffizienten werden dann als gleitender Mittelwert aktualisiert, um
Instabilitäten zu vermeiden. Die Autoren zeigen, dass diese Strategie sowohl
die Konvergenz beschleunigt als auch die Genauigkeit der Lösungen verbessert,
weil sie ein Ungleichgewicht zwischen physikalischen und Datengradienten
verhindert -- ein bekanntes Problem bei PINNs, das als \emph{gradient
pathology} bezeichnet wird.

% -----------------------------------------------------------------------
\subsection*{Experimente und Ergebnisse}

Die Methode wird an drei laminaren Benchmarks validiert:
\begin{itemize}
    \item \textbf{Kovasznay-Strömung (2D, stationär):} Relative $L_2$-Fehler
        im Bereich $10^{-5}$ bis $10^{-3}$; VV-Formulierung übertrifft VP
        bei großen Netzwerken.
    \item \textbf{Zylinderumströmung (2D, instationär, $Re = 100$):}
        Qualitativ und quantitativ gute Rekonstruktion der Wirbelablösung
        über einen vollen Ablösezyklus; VV erzielt durchgehend geringere
        Fehler als VP.
    \item \textbf{Beltrami-Strömung (3D, instationär):} Beide Formulierungen
        liefern genaue Lösungen; VV zeigt erneut bessere Genauigkeit, mit
        relativen Fehlern unter $0{,}3\,\%$ für alle Geschwindigkeitskomponenten.
\end{itemize}

Zusätzlich werden erste PINN-Simulationen turbulenter Kanalströmungen
bei $Re_\tau \approx 1000$ mit der VP-Formulierung und DNS-Referenzdaten
präsentiert. Die VP-NSFnet-Simulationen halten die Turbulenz über mehrere
konvektive Zeiteinheiten aufrecht, mit Geschwindigkeitsfehlern unter $10\,\%$,
wobei der Druckfehler bis zu $17\,\%$ erreichen kann.

% -----------------------------------------------------------------------
\subsection*{Bezug zur Masterarbeit}

Obwohl die NSFnets-Arbeit konzeptionell verwandt ist, unterscheidet sie sich in
mehreren wesentlichen Aspekten vom Ansatz der Masterarbeit -- diese Unterschiede
sind für die Positionierung der eigenen Arbeit besonders relevant:

\paragraph{Stokes- statt Navier-Stokes-Regime.}
Die Masterarbeit betrachtet die Sedimentation von Kristallen in Magma im
Stokes-Regime (sehr kleine Reynolds-Zahlen, $Re \ll 1$). In diesem Grenzfall
entfallen der nichtlineare Trägheitsterm $(\mathbf{u} \cdot \nabla)\mathbf{u}$
und der Zeitterm $\partial_t \mathbf{u}$, sodass die Stokes-Gleichungen
instantan eine elliptische PDE darstellen. Die bei NSFnets behandelten
Navier-Stokes-Gleichungen sind allgemeiner, aber für das geophysikalische
Problem nicht unmittelbar anwendbar. Das Stokes-Regime erlaubt außerdem die
Nutzung analytischer Referenzlösungen (z.\,B.\ Hadamard-Rybczynski) als
Ausgangspunkt für das Residuallernverfahren.

\paragraph{Surrogatmodell statt Solver.}
NSFnets verwenden PINNs als \emph{Solver}, der für jede neue Geometrie
neu trainiert werden muss. Die Masterarbeit verfolgt einen anderen Ansatz:
Ein UNet- bzw.\ FNO-basiertes Surrogatmodell wird einmalig auf einem
Datensatz von LaMEM-Simulationen trainiert und kann dann für neue
Kristallkonfigurationen \emph{ohne Neutraining} schnelle Vorhersagen liefern.
Die Generalisierungsfähigkeit über verschiedene Kristallanzahlen ($N$-zu-$N+m$)
stellt dabei den genuinen wissenschaftlichen Beitrag dar.

\paragraph{Strukturelle Erfüllung der Inkompressibilität.}
Die VV-Formulierung in NSFnets nutzt denselben mathematischen Mechanismus,
der auch in der Masterarbeit durch die \emph{Stromfunktionsformulierung}
realisiert wird: Statt die Divergenzfreiheitsbedingung $\nabla \cdot
\mathbf{u} = 0$ als Soft-Constraint in der Verlustfunktion zu erzwingen,
wird sie durch die Wahl der Ausgabegröße automatisch erfüllt. Für 2D-Stokes-
Strömungen gilt $\mathbf{u} = \nabla \times \psi$ mit der Stromfunktion
$\psi$, was $\nabla \cdot \mathbf{u} \equiv 0$ per Konstruktion garantiert.
NSFnets belegen, dass dieser Ansatz im Navier-Stokes-Kontext (VV-Formulierung)
tatsächlich zu besserer Genauigkeit führt als die explizite Soft-Constraint-
Variante -- ein wichtiges Argument, das die Designentscheidung der Masterarbeit
stützt.

\paragraph{Physik-informierte Verlustfunktionen.}
Die in NSFnets vorgestellte dynamische Gewichtungsstrategie ist direkt auf die
Masterarbeit übertragbar. Das dort implementierte adaptive Gewichtungsschema
zwischen Daten- und Physikloss basiert auf demselben Grundprinzip: Der
Gradientstatistik der verschiedenen Verlustterme wird verwendet, um ein
Ungleichgewicht zu verhindern. Die Experimente in NSFnets zeigen, dass ein zu
geringes Gewicht auf Randbedingungen ($\alpha = 1$) zu deutlich schlechteren
Ergebnissen führt -- eine Beobachtung, die mit den Erfahrungen in der
Masterarbeit konsistent ist.

\paragraph{Residuallernverfahren.}
Das Konzept des \emph{residual learning} -- also das Lernen von Korrekturen zu
einer bekannten Referenzlösung ($\mathbf{v}_\text{total} = \mathbf{v}_\text{Stokes}
+ \Delta\mathbf{v}_\text{gelernt}$) -- findet in NSFnets keine direkte
Entsprechung, da die Autoren vollständige Lösungen approximieren. Für das
Stokes-Regime der Masterarbeit bietet das Residuallernen jedoch einen klaren
Vorteil: Die analytische Einzelkristall-Lösung ist bekannt, und das Netzwerk
muss nur die Wechselwirkungskorrekturen zwischen Kristallen lernen, was den
Suchraum erheblich einschränkt und physikalisch sinnvollere Zwischenzustände
während des Trainings garantiert.

\end{document}